
où @@M0@@ est l'%
anneau @@M1@@, filtré par le poids par rapport aux @@M2@@, @@M3@@ étant de poids @@M4@@, et pour l'élément @@M5@@.

On plonge @@M6@@ dans @@M7@@ par les fonctions symétriques, et on obtient ainsi un homomorphisme injectif : @@M8@@~; il vient alors :
@@M9@@
On en déduit, d'après (6.2.1) :
@@M10@@
et on en tire :
@@M11@@
Enfin,
@@M12@@
Or, d'après (6.5.3), @@M13@@
soit : @@M14@@.

\textbf{Corollaire 6.6.2.} \uline{La relation} 6.6.1.iii) \uline{est vraie pour tout élément} @@M15@@ \uline{d'augmentation} @@M16@@.

Pour le montrer nous utiliserons le lemme :

\textbf{Lemme 6.6.3.} \uline{Pour tout} @@M17@@, @@M18@@ \uline{est un} @@M19@@-\uline{idéal de} @@M20@@.

On peut supposer @@M21@@. Soit @@M22@@~; on peut écrire @@M23@@ les @@M24@@ étant les composantes isobares de poids @@M25@@ des @@M26@@ (3.6). Pour tout @@M27@@, les polynômes @@M28@@ sont donc des polynômes par rapport aux @@M29@@ premiers coefficients de @@M30@@. Si @@M31@@, ils sont donc nuls, et par suite, les @@M32@@ appartiennent à @@M33@@, et 
